\begin{satz}[Skaleninvariantes Sinus-Integral auf Primzahlvierlingen]
	Sei
	\[
	(E,a,b,c)=(p,p+2,p+6,p+8)
	\]
	ein Primzahlvierling. Für jedes
	\[
	q\in\{E,a,b,c\}
	\]
	definiere die Funktion
	\[
	f_q:[0,q]\to\mathbb{R},\qquad
	f_q(x)=\frac{1}{q}\sin\!\left(\frac{\pi x}{q}\right).
	\]
	Dann gilt
	\[
	\int_0^E f_E(x)\,dx
	=
	\int_0^a f_a(x)\,dx
	=
	\int_0^b f_b(x)\,dx
	=
	\int_0^c f_c(x)\,dx
	=
	\frac{2}{\pi}.
	\]
	Außerdem gilt für jedes \(q\in\{E,a,b,c\}\) die Spiegelsymmetrie
	\[
	\int_0^q f_q(q-x)\,dx=\int_0^q f_q(x)\,dx.
	\]
\end{satz}

\begin{beweis}
	Fixiere zunächst ein beliebiges \(q>0\). Wir betrachten
	\[
	f_q(x)=\frac{1}{q}\sin\!\left(\frac{\pi x}{q}\right).
	\]
	
	\textbf{Schritt 1: Bestimmung einer Stammfunktion.}
	Setze
	\[
	F_q(x):=-\frac{1}{\pi}\cos\!\left(\frac{\pi x}{q}\right).
	\]
	Dann folgt mit der Kettenregel
	\[
	F_q'(x)
	=
	-\frac{1}{\pi}\left(-\sin\!\left(\frac{\pi x}{q}\right)\right)\frac{\pi}{q}
	=
	\frac{1}{q}\sin\!\left(\frac{\pi x}{q}\right)
	=
	f_q(x).
	\]
	Also ist \(F_q\) eine Stammfunktion von \(f_q\).
	
	\textbf{Schritt 2: Auswertung des Integrals.}
	Nach dem Hauptsatz der Differential- und Integralrechnung gilt
	\[
	\int_0^q f_q(x)\,dx
	=
	\left[
	-\frac{1}{\pi}\cos\!\left(\frac{\pi x}{q}\right)
	\right]_0^q.
	\]
	Einsetzen der Grenzen liefert
	\[
	\int_0^q f_q(x)\,dx
	=
	-\frac{1}{\pi}\cos(\pi)+\frac{1}{\pi}\cos(0)
	=
	-\frac{1}{\pi}(-1)+\frac{1}{\pi}(1)
	=
	\frac{2}{\pi}.
	\]
	Da dies für jedes \(q>0\) gilt, folgt insbesondere für
	\[
	q\in\{E,a,b,c\}
	\]
	die Gleichheit
	\[
	\int_0^E f_E(x)\,dx
	=
	\int_0^a f_a(x)\,dx
	=
	\int_0^b f_b(x)\,dx
	=
	\int_0^c f_c(x)\,dx
	=
	\frac{2}{\pi}.
	\]
	
	\textbf{Schritt 3: Spiegelungssymmetrie.}
	Nun berechnen wir
	\[
	f_q(q-x)
	=
	\frac{1}{q}\sin\!\left(\frac{\pi(q-x)}{q}\right)
	=
	\frac{1}{q}\sin\!\left(\pi-\frac{\pi x}{q}\right).
	\]
	Wegen der elementaren Identität
	\[
	\sin(\pi-u)=\sin(u)
	\]
	folgt
	\[
	f_q(q-x)=\frac{1}{q}\sin\!\left(\frac{\pi x}{q}\right)=f_q(x).
	\]
	Daher ist \(f_q\) bezüglich der Mitte \(x=q/2\) spiegelsymmetrisch. Insbesondere erhält man
	\[
	\int_0^q f_q(q-x)\,dx=\int_0^q f_q(x)\,dx.
	\]
	Damit ist alles bewiesen.
\end{beweis}

\begin{korollar}
	Für den ersten Primzahlvierling
	\[
	(E,a,b,c)=(5,7,11,13)
	\]
	gilt
	\[
	\int_0^5 \frac{1}{5}\sin\!\left(\frac{\pi x}{5}\right)\,dx
	=
	\int_0^7 \frac{1}{7}\sin\!\left(\frac{\pi x}{7}\right)\,dx
	=
	\int_0^{11} \frac{1}{11}\sin\!\left(\frac{\pi x}{11}\right)\,dx
	=
	\int_0^{13} \frac{1}{13}\sin\!\left(\frac{\pi x}{13}\right)\,dx
	=
	\frac{2}{\pi}.
	\]
	Numerisch ist
	\[
	\frac{2}{\pi}\approx 0.63661977236758.
	\]
\end{korollar}

\begin{bemerkung}
	Die Gleichheit der Integrale ist keine spezifische Folge der Primzahleigenschaft, sondern folgt aus der skaleninvarianten Form der Funktionen
	\[
	f_q(x)=\frac{1}{q}\,\phi\!\left(\frac{x}{q}\right)
	\]
	mit
	\[
	\phi(u)=\sin(\pi u).
	\]
	Allgemein gilt nämlich für jede integrierbare Funktion \(\phi:[0,1]\to\mathbb{R}\):
	\[
	\int_0^q \frac{1}{q}\phi\!\left(\frac{x}{q}\right)\,dx
	=
	\int_0^1 \phi(u)\,du.
	\]
	Der hier behandelte Sinusfall ist also ein spezielles Beispiel eines allgemeineren Skalenprinzips.
\end{bemerkung}